
\section {Summary of changes with respect to the {\em VALOR} Run 1+2+3  $\nu_{\mu}$-disappearance analysis }
\label{sec:changes}

This analysis is primarily a statistics update of our Run 1+2+3 $\nu_{\mu}$-disappearance analysis,
although certain analysis inputs and the fitting methodology have also been updated.
The following changes have been implemented with respect to our Run 1+2+3 analysis described in \cite{OurRun123DisappAnalysis3Flavour}:

\begin{itemize}

\item
Added the Run 4 SuperK data. Run 4 adds 3.5600$\times$10$^{20}$ POT (61 observed events) to the Run 1+2+3 exposure
of 3.010$\times$10$^{20}$ POT (59 observed events). The total Run 1+2+3+4 exposure is 6.5700$\times$10$^{20}$ POT
(\NskObserved observed events).

\item
The current analysis uses the most recent official SuperK MC samples (2013a) 
generated using the 11btunedv3.2 JNUBEAM flux.
All MC templates, systematic parameter response functions (cubic splines)
and systematic parameter covariance matrices have been regenerated with the new MC samples.

\item
The following 4 MC templates have been added for completeness.
{
\begin{itemize}
\item  oscillated $\nu_{e} \rightarrow \nu_{\mu}$ CCQE,
\item  oscillated $\nu_{e} \rightarrow \nu_{\mu}$ CC1$\pi$,
\item  oscillated $\nu_{e} \rightarrow \nu_{\mu}$ CC coherent and
\item  oscillated $\nu_{e} \rightarrow \nu_{\mu}$ CC other
\end{itemize}
}
Since a separate official oscillated $\nu_{e} \rightarrow \nu_{\mu}$ SuperK sample is not available,
we emulated it by reweighting the $\nu_\mu$ SuperK MC sample with the $\nu_e$ flux
(the oscillated $\nu_{e} \rightarrow \nu_{\mu}$ are not statistically independent of the $\nu_\mu$ templates).
In the oscillation fit, the contents of those 4 newly-added MC templates are weighted 
with P($\nu_{e}\rightarrow\nu_{\mu}$).
The contribution of these additional templates at the typical oscillation point used in Sec.\ \ref{sec:spectrum_prediction}
is $\approx$0.13\%.

\item
SuperK events with reconstructed energy $>$30 GeV are now included in the 10-30 GeV bin after discussion with the T2K-SK group 
(the reconstructed energy binning is unchanged with respect to the Run 1+2+3 analysis).
These events appear infrequently in MC;
only 14 out of 448463 events (0.003\%) after 1 $\mu$-like ring selections 
(the highest proportion occur in the $\nu_e$ CCQE \& $\nu_e$ CC other
modes, 1.0\% and 0.5\% respectively).
% 1 out of 38281 numu CC1pi
% 1 out of 1766 numubar NC oth
% 6 out of 1218 nue CCQE
% 2 out of 201 nue CC oth
% 1 out of 3824 nue NC oth
% 3 out of 1343 osc nue
In the Run 1+2+3+4 dataset, the highest energy 1 $\mu$-like ring event is 6.99 GeV.

\item 
The BANFF fit has been updated.
This analysis uses version 5 of the Run 1+2+3+4 BANFF 2013 fit \cite{BANFF2013} (as does the other $\nu_{\mu}$-disappearance 
analysis \cite{KDIRun1234DisappAnalysis}).
Besides adding the Run 4 ND280 data, the BANFF fit method was also updated.
For the 2013 analyses, the ND280 $\nu_{\mu}$ group updated the division of reconstructed ND280 CC events that are used as 
inputs to the BANFF fit.
These events are now divided into CC0$\pi$-like, CC1$\pi$-like and CC other samples,
whereas for the Run 1+2+3 analysis they were divided into CCQE-like and CCnonQE samples.

\item 
The SuperK detector error matrix has been recreated using more SuperK atmospheric data \cite{SuperkErrors2013}.

\item 
  The 7 SuperK detector efficiency uncertainties and the 6 FSI+SI uncertainties have been combined into a single 6-parameter
  covariance matrix\footnote{
    In \cite{OurRun123DisappAnalysis3Flavour}, the 7 SuperK efficiency parameters included 4 parameters for $\nu_\mu$+$\bar{\nu_\mu}$ CCQE;
    3 energy-dependent parameters for ring counting efficiencies, and 1 energy-independent parameter for all other efficiencies.
    In this analysis, these 4 parameters are combined giving a total of 6 SuperK detector efficiency uncertainties,
    which are then combined with the 6 FSI+SI uncertainties to create the single 6-parameter
    covariance matrix.
  }.
  This reduces the number of nuisance parameters in the fit, allowing the fit to run faster.
  This method is used by all the other official T2K analyses, both $\nu_\mu$ disappearance and $\nu_e$ appearance.

\item 
The response functions (cubic splines) for the uncorrelated cross-section systematic parameters have been regenerated using the
following oscillation parameters: $\sin^2(2\theta_{12})$ = 0.8495, $\sin^2(2\theta_{13})$ = 0.098, $\sin^2(2\theta_{23})$ = 1.0, 
$\Delta m^2_{21}$ = 7.5 $\times$ 10$^{-5}$ eV$^{2}$/c$^{4}$, $|\Delta m^2_{32}|$ = 2.44 $\times$ 10$^{-3}$ eV$^{2}$/c$^{4}$, 
$\delta_{CP}$ = 0.0 and the normal mass hierarchy.

\item 
The SuperK energy scale systematic has changed from 2.3\% to 2.4\% \cite{SuperkEScale2013}.

\item 
The CC $\nu/\bar{\nu}$ cross section normalisation systematic has changed from 40\% to 20\% \cite{NIWGNuBarErr}.

\item 
During the review process of the T2K Run 1+2+3 $\nu_{\mu}$-disappearance PRL \cite{T2KRun123Disapp}, it was discovered that the 
choice of octant ($\theta_{23} \leq \pi/4$ or $\theta_{23} \geq \pi/4$) has a large effect on the confidence contour in 
$\sin^2(2\theta_{23})$. Therefore, we now fit $\sin^2\theta_{23}$ by default (though the ability to fit $\sin^2(2\theta_{23})$ 
has been retained).

\item 
The neutrino mass hierarchy assumption also has a large effect on the confidence contour in $|\Delta m^2_{32}|$ or $|\Delta m^2_{31}|$, 
and we therefore produce two results: one assumes the normal hierarchy and fits $|\Delta m^2_{32}|$, and the other assumes the inverted 
hierarchy and fits $|\Delta m^2_{31}|$.
% The inputs to the calculation of the oscillation probabilities are +$|\Delta m^{2}_{32}|$ for the normal
% hierarchy and -$|\Delta m^{2}_{31}|$ for the inverted hierarchy. 

\item 
In this analysis, all four \emph{non-23-sector} oscillation parameters are allowed to float in the fit, whereas they were fixed in our 
Run 1+2+3 analysis. 
We use constraints from the the 2012 PDG for $\Delta m^2_{21}$, $\sin^2\theta_{13}$ and $\sin^2\theta_{12}$ (see Tab.\ \ref{tab:pdg2012non23}), 
while $\delta_{CP}$ is allowed to float freely in the range [-$\pi$, +$\pi$].
When generating toy MC experiments, values of $\Delta m^2_{21}$, $\sin^2\theta_{13}$ and $\sin^2\theta_{12}$ are
thrown from a normal distribution with mean and standard deviation taken from Tab.\ \ref{tab:pdg2012non23},
whereas $\delta_{CP}$ is thrown uniformly within its allowed range. 

\item
  In this analysis, the spectral function systematic is restricted in the fit to the range [0,1], with a Gaussian constraint that
  favoured 0 (relativistic Fermi gas).
  This was also done in our Run 1+2+3 analysis, but, in that analysis,
  the spectral function systematic was thrown uniformly within its allowed range when generating toy MC experiments.
  In this analysis, toy MC experiments are generated using a Gaussian constraint, centered at 0, and bounded in the range [0,1],
  as per the NIWG recommendation.
  
\item
We have added 1D confidence intervals in $\sin^2\theta_{23}$ and $|\Delta m^2_{32}|$ (normal hierarchy) or $|\Delta m^2_{31}|$ 
(inverted hierarchy) using a new procedure that is an extension of the Feldman-Cousins method.
\end{itemize}


